\chapter{Opioid safety}
\index{Opioid safety}

\section*{Definition and Overview}
Opioid safety encompasses the clinical and self-care practices designed to minimise the risks associated with the prescription, dispensing, administration, and storage of opioid medications. While opioids are highly effective analgesics for managing moderate to severe pain, their narrow therapeutic window and high potential for abuse, physical dependence, and fatal respiratory depression make safety protocols critical. Implementing opioid safety standards involves clinical screening, strict dosing controls, monitoring for adverse events, educating patients on the risks of drug interactions (particularly with other central nervous system depressants), and establishing protocols for the safe storage and disposal of these potent substances.

\section*{Epidemiology}
Prescription opioid-related harm is a global public-health issue. Risk varies between countries and communities and is increased by high doses, long-acting or highly potent opioids, loss of tolerance, use with alcohol or sedatives, sleep-disordered breathing, older age, and some medical conditions. People with chronic pain, mental-health conditions, or a history of substance use may need additional support, but anyone taking an opioid can experience serious adverse effects.

\section*{Aetiology and Pathophysiology}
Opioid safety concerns arise from the specific pharmacological effects of opioids on the central nervous system and the physiological adaptations that occur during prolonged exposure.
\begin{itemize}
    \item \textbf{Respiratory depression pathway:} Opioids bind to mu-opioid receptors in the pre-B\"otzinger complex of the brainstem, suppressing the automatic respiratory rhythm and decreasing the ventilatory response to hypercapnia and hypoxia.
    \item \textbf{Synergistic CNS depression:} Co-administration of opioids with other GABA receptor agonists (such as benzodiazepines or alcohol) causes profound, synergistic depression of the central nervous system, leading to rapid loss of consciousness and respiratory arrest.
    \item \textbf{Tolerance development:} Chronic receptor stimulation leads to receptor desensitisation and downregulation, requiring higher doses to achieve the same analgesic effect and narrowing the margin between therapeutic and lethal doses.
    \item \textbf{Physical dependence:} Cellular adaptations in the locus coeruleus lead to severe physiological withdrawal symptoms if the opioid is abruptly ceased or reduced.
    \item \textbf{Cognitive impairment and motor coordination deficits:} Sedative effects impair psychomotor performance, increasing the risk of mechanical falls, motor vehicle accidents, and cognitive decline, particularly in older patients.
\end{itemize}

\section*{Clinical Features}
Features of opioid safety compromises can present as signs of chronic misuse, adverse drug events, or acute toxicity.
\begin{itemize}
    \item \textbf{Excessive sedation:} Unusually deep sleep, persistent drowsiness, or difficulty remaining awake during normal activities.
    \item \textbf{Cognitive slowing:} Confusion, memory impairment, or lack of coordination and balance.
    \item \textbf{Respiratory compromise:} Slow, shallow breathing, snoring or snorting sounds during sleep, and periods of apnoea.
    \item \textbf{Signs of misuse:} Requesting early refills, obtaining prescriptions from multiple doctors (``doctor shopping''), or using medications for reasons other than pain (such as anxiety or sleep).
    \item \textbf{Physical signs of dependence:} Sweating, tremors, dilated pupils, and abdominal cramping when a dose is delayed.
    \item \textbf{Adverse effects:} Severe constipation, urinary retention, dry mouth, and chronic itching.
\end{itemize}

\section*{Differential Diagnosis}
Discerning whether symptoms of somnolence or respiratory distress are due to opioid toxicity, drug interactions, or other medical issues is critical.
\begin{itemize}
    \item \textbf{Non-opioid drug toxicity:} Sedative-hypnotic overdose (such as benzodiazepines, z-drugs, or barbiturates) which can present with similar levels of obtundation but generally lack pinpoint pupils.
    \item \textbf{Carbon dioxide narcosis:} Respiratory failure due to severe chronic obstructive pulmonary disease (COPD) or obesity hypoventilation syndrome, presenting with somnolence and respiratory acidosis.
    \item \textbf{Sepsis or systemic infection:} Can cause altered mental state and lethargy, especially in older adults, but is accompanied by fever, tachycardia, and inflammatory markers.
    \item \textbf{Acute neurological events:} Ischaemic or haemorrhagic stroke, or post-ictal states, which present with focal neurological deficits or altered consciousness.
\end{itemize}

\section*{When to Seek Medical Care}
Patients or caregivers should contact their prescriber if they notice signs of worsening side effects, increased confusion, or if the prescribed dose is no longer controlling pain. Immediate emergency care is required if there are signs of opioid toxicity or respiratory compromise.

\textbf{Contact your local emergency services for:}
\begin{itemize}
    \item Inability to rouse the person, or responsiveness only to painful stimuli
    \item Respiratory rate of fewer than 8 breaths per minute, or long pauses between breaths
    \item Pale or blue-tinged skin, lips, or fingernails
    \item Pinpoint pupils in an unresponsive or barely conscious individual
    \item Choking, snorting, or gurgling sounds in someone who cannot be easily awakened
\end{itemize}

\section*{Diagnosis and Investigations}
Safety monitoring involves establishing baseline risk and actively screening for signs of toxicity or misuse.
\begin{itemize}
    \item \textbf{Opioid Risk Tool (ORT):} A validated clinical questionnaire completed prior to initiating therapy to assess the patient's risk of opioid misuse.
    \item \textbf{Prescription monitoring:} Where available and permitted, clinicians may review dispensing records for high-risk medicines as one part of a broader assessment; monitoring systems and privacy rules differ between countries.
    \item \textbf{Urine drug testing (UDT):} Periodic screening to verify compliance with the prescribed regimen and detect non-prescribed substances.
    \item \textbf{Pulse oximetry and capnography:} Monitoring blood oxygen saturation and carbon dioxide levels in hospitalised patients receiving intravenous opioids.
    \item \textbf{Hepatic and renal function tests:} Essential to assess the body's ability to metabolise and clear active metabolites of opioids, preventing drug accumulation and toxicity.
\end{itemize}

\section*{Management}
Managing opioid safety involves rigorous prescribing guidelines, active risk mitigation, and provision of emergency reversal agents.
\begin{itemize}
    \item \textbf{Universal precautions in prescribing:} Initiating therapy at the lowest dose, utilising short-acting formulations first, and establishing clear boundaries for refills.
    \item \textbf{Naloxone distribution (Take-home Naloxone):} Providing patients at high risk of overdose, or their families, with intranasal or intramuscular naloxone kits for emergency reversal.
    \item \textbf{Avoidance of co-prescribing:} Strictly limiting the concurrent prescription of opioids with benzodiazepines or other CNS depressants unless clinically unavoidable.
    \item \textbf{Structured tapering:} If a reduction is appropriate, agree a gradual, individualised plan with the prescriber. The rate depends on duration of use, dose, withdrawal symptoms, pain, and patient preference; abrupt discontinuation should generally be avoided unless there is an immediate safety emergency.
    \item \textbf{Multimodal analgesia:} Integrating non-opioid medications (paracetamol, NSAIDs) and non-pharmacological interventions (physiotherapy, psychological therapies) to minimise opioid requirements.
\end{itemize}

\section*{Complications}
\begin{itemize}
    \item Accidental or intentional overdose leading to fatal respiratory arrest
    \item Hypoxic brain injury or permanent neurological deficits secondary to prolonged respiratory depression
    \item Development of severe opioid use disorder, dependency, or illicit substance transition
    \item Increased risk of aspiration pneumonia due to depressed cough reflex and altered consciousness
    \item Hyperalgesia, where the pain becomes worse and more widespread over time
    \item Immunosuppression and endocrine deficiency with long-term use, increasing susceptibility to infections and osteoporosis
\end{itemize}

\section*{Prognosis and Outcomes}
Risk can be reduced through shared decisions, regular review, careful dosing, avoidance of dangerous combinations, and access to naloxone where available. There is no single dose threshold that is safe for everyone; risk depends on the person, opioid, duration, co-medications, and clinical setting. If long-term treatment is reduced, a slow, supportive, individualised plan is safer than abrupt stopping.

\section*{Prevention and Self-Care}
\begin{itemize}
    \item Never share opioid medicines with others or use another person's prescribed pain medication.
    \item Store all opioid medications in a secure, locked container, completely out of reach of children, adolescents, and visitors.
    \item Avoid alcohol entirely while taking opioids, and discuss all other medications—especially sleep or anxiety drugs—with your doctor or pharmacist.
    \item Dispose of unused or expired opioid medications safely by returning them to a local community pharmacy for secure destruction.
    \item Never crush, chew, or break extended-release opioid tablets, as this can cause a rapid, fatal release of the entire dose (``dose dumping'').
\end{itemize}

\section*{Sources and Further Reading}
The following organisations provide reliable resources and guidelines on opioid safety and prescription management:
\begin{itemize}
    \item World Health Organization. \textit{Opioid overdose.} https://www.who.int/news-room/fact-sheets/detail/opioid-overdose
    \item World Health Organization. \textit{Community management of opioid overdose.} https://www.who.int/publications-detail-redirect/9789241548816
    \item Centers for Disease Control and Prevention. \textit{Clinical practice guideline for prescribing opioids for pain.} https://www.cdc.gov/overdose-prevention/hcp/clinical-guidance/
\end{itemize}
